\section*{\fontsize{18pt}{18pt}\selectfont\MakeUppercase{\de}\\TÓM TẮT}

Trong bối cảnh các trường đại học ngày càng mở rộng quy mô đào tạo, số lượng người dùng, thiết bị đầu cuối và dịch vụ trực tuyến trong khuôn viên tăng nhanh, việc xây dựng một hệ thống mạng Campus có khả năng quản lý tập trung, vận hành ổn định và mở rộng linh hoạt trở thành yêu cầu quan trọng. Các mô hình mạng truyền thống thường phụ thuộc nhiều vào cấu hình thủ công trên từng thiết bị, gây khó khăn khi cần bổ sung khoa, phòng học, phòng máy hoặc kết nối thêm cơ sở phụ.

Xuất phát từ thực tế đó, đề tài tập trung nghiên cứu và xây dựng mô hình mạng Campus Network sử dụng Software-Defined Networking (SDN) và Software-Defined Wide Area Networking (SD-WAN). Hệ thống được định hướng theo kiến trúc Core--Distribution--Access, phân chia VLAN theo từng khu vực chức năng, tổ chức địa chỉ IP có cấu trúc và áp dụng các cơ chế định tuyến, bảo mật, dự phòng nhằm bảo đảm khả năng kết nối cho hoạt động đào tạo, nghiên cứu và quản trị trong trường đại học.

Trong mô hình đề xuất, SDN Controller đóng vai trò quản lý tập trung các thiết bị chuyển mạch, hỗ trợ cấu hình chính sách, giám sát luồng dữ liệu và đơn giản hóa quá trình vận hành mạng nội bộ Campus. Bên cạnh đó, SD-WAN được sử dụng để kết nối Campus chính với các cơ sở phụ thông qua nhiều đường truyền WAN, giúp tăng tính linh hoạt, khả năng dự phòng và tối ưu lưu lượng giữa các địa điểm. Các kỹ thuật nền tảng như VLAN, DHCP, OSPF, STP/RSTP, ACL và dự phòng liên kết được tích hợp để hoàn thiện mô hình triển khai.

\newpage

Nội dung báo cáo được tổ chức thành năm chương chính như sau:

\begin{itemize}
    \item[-] \textbf{Chương 1 -- Tổng quan đề tài:} Trình bày lý do chọn đề tài, mục tiêu, phạm vi, phương pháp thực hiện và bố cục báo cáo.
    \item[-] \textbf{Chương 2 -- Cơ sở lý thuyết:} Trình bày các nội dung về Campus Network, mô hình phân cấp, VLAN, định tuyến, SDN và SD-WAN.
    \item[-] \textbf{Chương 3 -- Mô hình mạng đề xuất:} Mô tả topology, phân hoạch VLAN, địa chỉ IP và chính sách quản lý.
    \item[-] \textbf{Chương 4 -- Triển khai và kiểm thử:} Trình bày phương án triển khai, cấu hình các dịch vụ mạng và các kịch bản kiểm thử.
    \item[-] \textbf{Chương 5 -- Kết luận và hướng phát triển:} Tổng kết kết quả đạt được, đánh giá hạn chế và đề xuất hướng phát triển.
\end{itemize}

Kết quả kỳ vọng của đề tài là xây dựng được một khung thiết kế mạng Campus có tính thực tiễn, dễ quản trị, dễ mở rộng và phù hợp với môi trường đào tạo nhiều khu vực chức năng. Mô hình này giúp giảm phụ thuộc vào cấu hình thủ công, hỗ trợ quản lý chính sách tập trung, nâng cao khả năng sẵn sàng của hệ thống và tạo nền tảng để tiếp tục phát triển các chức năng giám sát, tự động hóa và tối ưu lưu lượng trong tương lai.

\newpage
